\documentclass[11pt,letterpaper]{article}

%% ── Packages ──────────────────────────────────────────────────────────────────
\usepackage[T1]{fontenc}
\usepackage[utf8]{inputenc}
\usepackage{lmodern}
\usepackage[margin=1in]{geometry}
\usepackage{microtype}
\usepackage{amsmath,amssymb}
\usepackage{booktabs}
\usepackage{array}
\usepackage{tabularx}
\usepackage{multirow}
\usepackage{graphicx}
\usepackage{xcolor}
\usepackage[colorlinks=true,linkcolor=blue!60!black,citecolor=blue!60!black,urlcolor=blue!60!black]{hyperref}
\usepackage{setspace}
\usepackage{parskip}
\usepackage{titlesec}
\usepackage{enumitem}
\usepackage{caption}
\usepackage{url}
\usepackage{fancybox}
\usepackage{float}

%% ── Custom commands ───────────────────────────────────────────────────────────
\newcommand{\doi}[1]{\href{https://doi.org/#1}{\texttt{doi:#1}}}

%% ── Typography ────────────────────────────────────────────────────────────────
\setstretch{1.10}
\titleformat{\section}{\normalfont\large\bfseries}{\thesection}{0.7em}{}
\titleformat{\subsection}{\normalfont\normalsize\bfseries}{\thesubsection}{0.6em}{}
\setlength{\parindent}{0pt}
\setlength{\parskip}{6pt}

%% ── Metadata ─────────────────────────────────────────────────────────────────
\hypersetup{
  pdftitle  = {A Governed QC-AI-HPC Information-Systems Architecture for Quantum Simulation of Non-druggable NSCLC Tumor Suppressor Mutations},
  pdfauthor = {Doron Cohen},
}

%% ══════════════════════════════════════════════════════════════════════════════
\begin{document}

%% ── Title block ──────────────────────────────────────────────────────────────
\begin{center}
{\Large\bfseries A Governed QC--AI--HPC Information-Systems Architecture\\
for Quantum Simulation of Non-druggable NSCLC Tumor Suppressor Mutations}\\[10pt]
{\normalsize Doron Cohen}\\
{\small PhD/DTech Candidate, Center for Information Systems and Technology (CISAT)\\
Claremont Graduate University (CGU)\\
\href{mailto:doron.cohen@cgu.edu}{doron.cohen@cgu.edu}\\
Supervisor: Prof.\ Itamar Shabtai}\\[6pt]
{\small May 2026 \quad|\quad Intended arXiv category: cs.ET}\\
{\small Author-selected arXiv license: CC BY-NC-ND 4.0}
\end{center}

\noindent\rule{\textwidth}{0.4pt}

%% ── Abstract ─────────────────────────────────────────────────────────────────
\begin{center}\textbf{Abstract}\end{center}

% CORRECTION #1 / #5: Assertive IS framing replaces defensive opener;
% all clinical/chemistry disclaimers moved to Limitations (Section 10).
Regulated quantum-oncology simulation is an information-systems problem before it is a
computational-chemistry problem.  Governing \emph{when} to escalate a workload to a quantum
kernel---and producing an auditable record of that decision---requires an orchestration
architecture, not merely a quantum algorithm.  The contribution of this paper is precisely
that architecture: a governed QC--AI--HPC orchestration system that coordinates regulated
molecular inputs, active-space classification, selective computational escalation,
backend-specific compilation, AI/HPC support services, stakeholder-facing visualisation,
and Part-11 provenance for probabilistic quantum simulation outputs.

The architecture has three layers.  Layer~1 is the quantum-kernel layer, anchored by SqDRIFT
on IBM Heron~r3 for frontier fermionic-Hamiltonian workloads.  Layer~2 is the hybrid-substrate
layer, covering classical preprocessing and postprocessing, GPU/HPC services, AI calibration,
and emerging AI-decoded quantum-error-correction substrates including NVIDIA CUDA-Q, NVQLink,
and Ising-class learned decoders.  Layer~3 is the governed orchestration and provenance layer,
which routes tasks across the first two layers while emitting a nine-element P1--P9 provenance
record---the paper's core information-systems contribution---at each stage of execution.
A design-science evaluation protocol is proposed combining technical harnesses, comparative
benchmarking, expert review, and case-based walkthroughs.  The current implementation executes
on a PennyLane noiseless simulator (Phase~3A) using physically motivated proxy Hamiltonians
for pipeline validation; IBM Heron~r3 hardware execution with PySCF-derived molecular
Hamiltonians is Phase~3B.  No energy result in this paper should be interpreted as a
quantum-computed molecular binding energy.

\smallskip
\noindent\textbf{Keywords:} information systems; design science research; emerging computing;
QC--AI--HPC orchestration; quantum computing; high-performance computing; artificial intelligence;
precision oncology; TP53; NSCLC; SqDRIFT; IBM Heron; NVIDIA NVQLink; FDA~21~CFR~Part~11;
provenance.

\vfill
\noindent\rule{\textwidth}{0.4pt}
{\small\textit{\copyright{} 2026 Doron Cohen.  Licensed under CC BY-NC-ND~4.0.}}

\clearpage
%% ── §1 Introduction ──────────────────────────────────────────────────────────
\section{Introduction}

% CORRECTIONS #1 / #5: Intro rewritten to lead with IS architecture claim.
% One "rules of engagement" paragraph replaces scattered defensive disclaimers.
The central thesis of the dissertation proposal underlying this article is direct: the
contribution is the \textbf{architecture}; simulation is the application; the patient case is
the motivation~\cite{cohen2026proposal}.  This paper asks a systems question---what
architecture is required so that a quantum-enabled oncology simulation workflow can be
executed, reviewed, rerun, and governed without pretending that all computation occurs in
a single deterministic environment?---and answers it with a governed orchestration layer
implementing a real-time audit log for hybrid quantum--classical--AI pipelines.

\textbf{Scope of this paper.}  The architecture specifically addresses the information-systems
challenge of regulated computational escalation: deciding when, why, and under what recorded
conditions work moves from classical and AI services to a quantum kernel.  The motivating domain
is non-druggable NSCLC tumor-suppressor mutation simulation; the stress-test case is
TP53 p.Cys275Phe (C275F).  NSCLC accounts for approximately 2.1~million new diagnoses and
1.53~million deaths annually (GLOBOCAN~2022), and approximately 60\% of patients carry
mutations for which no approved targeted therapy exists.  Clinical validation, novel
binding-energy estimates, and therapeutic recommendations are outside the scope of this
architecture paper; these boundaries are design choices, not deficits.  Readers seeking the
chemistry or clinical evidence base should consult the literature cited in the Limitations
section.

Three architectural challenges motivate the design.  First, a regulated oncology simulation
workflow must \emph{govern escalation}: the decision rule routing work from classical and AI
methods to quantum kernels must be explicit, revisable, and recorded in a compliance-grade
audit trail.  Second, the workflow must manage a bifurcating hardware landscape in which
IBM Heron/Qiskit and NVIDIA CUDA-Q/NVQLink paths expose different calibration records,
compiler behaviours, and audit surfaces.  Third, quantum simulation outputs are
probabilistic distributions, not deterministic values, so the compliance record must
operationalise reproducibility as statistical equivalence under recorded conditions rather
than bitwise replay.

This paper therefore presents a bifurcation-aware QC--AI--HPC orchestration architecture.
It treats IBM Heron~r3 and SqDRIFT as the dissertation's quantum-kernel anchor; treats NVIDIA
CUDA-Q, NVQLink, and Ising-class decoders as an emerging hybrid-substrate pattern where
applicable; and locates the information-systems contribution in the orchestration layer that
makes those choices explicit, auditable, and revisable~\cite{heron2026ibm,ising2026nvidia,cudaqec2026}.

%% ── §2 Design-Science Framing ────────────────────────────────────────────────
\section{Design-Science Framing and Contributions}

The paper follows design science research (DSR): the research contribution is an artifact
designed and evaluated for utility, rigor, and novelty~\cite{hevner2004,peffers2007}.  The
artifact is a workflow architecture and provenance schema for a class of regulated
QC--AI--HPC simulation workflows, not a molecule, a clinical endpoint, or a standalone quantum
circuit.

Three dissertation-derived contributions are translated into arXiv form:

\begin{enumerate}[leftmargin=1.5em,itemsep=3pt]
\item \textbf{BQP-relevant classification as routing policy.}  A tumor-suppressor mutation
target becomes operationally BQP-relevant when the active-space boundary exceeds tractable
classical accuracy under available resources and the corresponding fermionic-Hamiltonian problem
admits a known quantum algorithmic formulation.  This is a decision rule for computational
escalation, not a formal hardness proof.

\item \textbf{Part-11-governed QC--AI--HPC orchestration.}  The architecture is designed for
oncology mutation simulation workflows in which regulated inputs, backend execution, AI/HPC
services, and stakeholder outputs are coordinated by a governable information system rather
than by ad hoc scripts.

\item \textbf{Provenance for probabilistic quantum outputs and learned decoders.}  A
nine-element P1--P9 record operationalises auditability for quantum simulation outputs.  P9
is included because emerging learned quantum-error-correction (QEC) decoders make model
lineage part of the computational record.
\end{enumerate}

The evaluation plan follows FEDS, which treats DSR evaluation as a deliberate strategy rather
than a single end-stage test~\cite{venable2016}.  The proposed evaluation therefore combines
technical record-completeness tests, comparative workflow benchmarks, expert review, and
case-based walkthroughs.

Table~\ref{tab:hevner} maps the three contributions against Hevner et al.'s seven DSR
guidelines~\cite{hevner2004} to make the design-science structure explicit for IS reviewers.

\begin{table}[H]
\centering
\caption{Three contributions mapped to Hevner's seven DSR guidelines~\cite{hevner2004}.}
\label{tab:hevner}
\small
\begin{tabularx}{\textwidth}{@{}lX@{}}
\toprule
\textbf{DSR Guideline} & \textbf{How this paper addresses it} \\
\midrule
Design as artifact &
The three-layer architecture, P1--P9 provenance schema, RBAC/JWT access layer, and working
prototype (Phase~3A live; Phase~3B planned) collectively constitute the artifact, publicly
accessible at \url{https://nsclc-quantum-viewer.netlify.app}. \\[4pt]
Problem relevance &
Section~3 establishes three domain gaps: clinical (no approved therapy for $\sim$60\% of
NSCLC patients), computational (classical FCI ceiling at $\sim$18 active electrons), and
systems (no governed orchestration connecting these layers exists). \\[4pt]
Design evaluation &
Section~8 specifies a four-layer evaluation protocol: technical harnesses ($\geq$95\%
completion, 100\% P1--P9 emission), comparative benchmarks against four alternatives,
expert-panel scoring (n\,=\,4--6, Hevner criteria, 1--5 scale), and case-based walkthrough
of two mutation scenarios. \\[4pt]
Research contributions &
Sections~2 and~6: BQP classification rule (novel operational escalation criterion), P1--P9
schema (first provenance framework for probabilistic quantum simulation outputs), RBAC/JWT
access layer (Part-11 access control and user attribution).  None are present in existing
orchestration alternatives. \\[4pt]
Research rigor &
DSR methodology follows Hevner et al.\ (2004)~\cite{hevner2004} and Peffers et al.\
(2007)~\cite{peffers2007}.  Evaluation strategy follows FEDS (Venable et al.,
2016)~\cite{venable2016}, combining ex~ante artificial evaluation with ex~post naturalistic
evaluation. \\[4pt]
Design as search process &
The three-layer architecture emerged from iterative response to three domain constraints
(Section~1): the escalation-governance requirement, the bifurcating hardware landscape, and
the probabilistic-output compliance problem. \\[4pt]
Communication of research &
arXiv submission targets cs.ET for emerging-computing and IS audiences.  The companion
prototype makes the artifact simultaneously inspectable by QC specialists, IS scholars,
regulatory readers, and clinical informatics reviewers. \\
\bottomrule
\end{tabularx}
\end{table}

%% ── §3 Motivating Use Case ────────────────────────────────────────────────────
\section{Motivating Use Case and Active-Space Requirements}

The motivating case is TP53~p.Cys275Phe (C275F) in NSCLC.  In the dissertation proposal,
C275F functions as a requirements-driving case: it creates a need for mutation-aware ingestion,
active-space definition, selective escalation, and stakeholder-facing explanation~\cite{cohen2026proposal}.

TP53 is selected because its mutation landscape is architecturally heterogeneous.  Different
mutations produce different active-space sizes, placing distinct demands on qubit count,
circuit depth, telemetry volume, and backend routing.  C275F is used because its active-space
boundary---approximately 24 active electrons at the local binding site---sits at the
classical--quantum frontier where escalation decisions are genuinely contested and the
governance architecture is most severely stressed.

Table~\ref{tab:tiers} expresses the requirements as systems-design tiers.  The active-space
estimates determine qubit demand, circuit depth, telemetry volume, backend routing, and
provenance granularity.

\begin{table}[H]
\centering
\caption{Dissertation-derived workload tiers used by the orchestration architecture.}
\label{tab:tiers}
\small
\begin{tabular}{@{}lccp{6.5cm}@{}}
\toprule
\textbf{Tier} & \textbf{Active e$^{-}$} & \textbf{Approx.\ qubits} & \textbf{Architectural role} \\
\midrule
Reference molecular benchmark & 32 & 100 &
  External SqDRIFT/Heron precedent (IBM \textit{Science}, March 2026) \\
C275F local active site & $\sim$24 & $\sim$48 &
  Phase 3A proof-of-concept; simulator validation; near-term quantum-kernel test \\
C275F full active site  & $\sim$44 & $\sim$88 &
  \textbf{Only target within today's hardware ceiling.} Merz et al.~\cite{merz2026protein} demonstrated 94 qubits on IBM Heron~r2 (May 2026, PDB:~2OCJ active-site verified); Phase~3B pending IBM Quantum Network access \\
KEAP1 local site (G333C) & $\sim$104 & $\sim$208 &
  Fault-tolerant QPU required ($\sim$2030+); PDB:~1U6D+2FLU coordinate-verified \\
KEAP1 full Nrf2 interface & $\sim$155 & $\sim$310 &
  Fault-tolerant QPU required; PDB:~2FLU coordinate-verified (19 interface residues) \\
STK11/LKB1 local site (D194) & $\sim$76 & $\sim$152 &
  Fault-tolerant QPU required; PDB:~2WTK coordinate-verified (D194A backbone position) \\
STK11/LKB1 full ATP pocket & $\sim$152 & $\sim$304 &
  Fault-tolerant QPU required; PDB:~2WTK coordinate-verified (8\,\AA{} shell) \\
\bottomrule
\end{tabular}
\end{table}

The local tier tests end-to-end execution: de-identified input handling, active-space
construction, fermion-to-qubit mapping, backend routing, error-mitigation metadata, AI/HPC
support, stakeholder visualisation, and P1--P9 provenance emission.  The full tier tests
whether the same architecture scales without changing its governance contract.

The platform addresses seven scientifically classified NSCLC mutation targets selected
on three criteria: (1)~clinical significance in NSCLC specifically, (2)~absence of an
approved targeted therapy, and (3)~BQP-relevant active space exceeding classical
tractability.  Hardware feasibility on current quantum computers was deliberately
\emph{not} a selection criterion---the architecture must be mutation-agnostic and
future-ready.  The seven targets are: TP53~C275F (Class~A), KEAP1~LOF, KEAP1~G333C,
KEAP1~R320Q (all Class~B), STK11~LKB1~(LOF), STK11~F354L, and STK11~D194N (all
Class~A).  STK11 and KEAP1 targets were selected for their clinical urgency in NSCLC:
STK11/LKB1 loss-of-function occurs in $\sim$20\% of NSCLC adenocarcinomas and is the
primary driver of immunotherapy (PD-1/PD-L1) resistance; KEAP1 mutations co-occur in
$\sim$10--15\% of NSCLC cases and together define the ``double-negative''
immunotherapy non-responder profile---one of the most urgent unsolved problems in
thoracic oncology.

Coordinate-level active-space analysis from PDB structures (2OCJ, 1U6D, 2FLU, 2WTK)
establishes a critical hardware-tier distinction.  TP53~C275F is the \emph{only} target
whose full active site ($\sim$44 electrons, $\sim$88 qubits; PDB:~2OCJ) falls within the
94-qubit ceiling independently demonstrated by Merz et al.~\cite{merz2026protein}.
The six KEAP1 and STK11 targets require 152--310 qubits (Table~\ref{tab:tiers}) and are
designated fault-tolerant QPU targets ($\sim$2030+).  Their inclusion in the platform
validates the architecture's mutation-agnostic routing and governance pipeline across the
full hardware-tier spectrum---demonstrating that the same provenance schema and compliance
framework apply whether the backend is a 2-qubit PennyLane simulator today or a
fault-tolerant QPU tomorrow.  C275F is the primary stress-test case because it is the
only target that is simultaneously quantum-necessary (Class~A---no classical path at
chemical accuracy) and executable on today's IBM Heron hardware.

\textbf{Translational chain and governance claim.}  The architecture produces
quantum-accurate ground-state energy differentials ($\Delta E$) between wild-type and
mutant active sites.  $\Delta E$ is not the terminal output; it is the necessary input to
downstream binding free energy estimation via QM/MM frameworks, where
$\Delta G_{\text{binding}} = \Delta E_{\text{total}} + \Delta G_{\text{solvation}} -
T\Delta S$.  Classical density-functional methods introduce systematic errors of
2--4~kcal/mol in correlated active sites; at this error level, predicted binding affinities
shift by a factor of five to six, sufficient to misclassify a viable drug candidate as
inactive.  Quantum-accurate $\Delta E$ reduces this error to approximately 0.5~kcal/mol,
within the threshold of chemical accuracy.  For the IS architecture, this chain establishes
the scope of the governance claim: the output the system audits and cryptographically seals
is the minimum computational signal required to justify downstream drug-design
investment---making the provenance record a prerequisite for translational use, not a
compliance formality.

%% ── §4 Three-Layer Architecture ──────────────────────────────────────────────
\section{Bifurcation-Aware Three-Layer Architecture}

\subsection{Architectural principle}

The dissertation's architectural rule is: hardware-agnostic at the workflow level;
vendor-specific at the compilation and error-mitigation layer~\cite{cohen2026proposal}.
A regulated workflow cannot pretend that superconducting, trapped-ion, neutral-atom, photonic,
simulator, GPU, and HPC backends expose the same latency, calibration records, compiler
behaviour, error channels, or audit data.  The system hides unnecessary complexity from users
but must not hide substrate-specific lineage from the compliance record.

The post-2026 hybrid-computing landscape strengthens this rule.  IBM Heron~r3 is a 156-qubit
Heron-family revision whose documented r3 improvements target coherence, gate fidelity, and
readout performance~\cite{heron2026ibm}.  NVIDIA's Ising and CUDA-Q QEC materials position
GPUs, NVQLink, CUDA-Q, and learned decoders as real-time hybrid-substrate components for
quantum calibration and QEC~\cite{ising2026nvidia,cudaqec2026}.  The architecture therefore
spans a vendor bifurcation: Qiskit Runtime and IBM Heron where the IBM path is selected;
CUDA-Q, NVQLink, and Ising-class services where a non-IBM or GPU-coupled substrate is
appropriate.

Figure~\ref{fig:arch} summarises the three-layer architecture.

\begin{figure}[H]
\centering
\small
\fbox{%
\begin{minipage}{0.92\textwidth}
\small
\noindent\textbf{Regulated inputs} $\rightarrow$
De-identified mutation case, structure/model evidence, active-space tier,
study objective, admissible backend policy

\medskip
\noindent\colorbox{gray!15}{\parbox{\dimexpr\textwidth-2\fboxsep}{%
\textbf{Layer~3: governed orchestration and provenance layer}\\
Selective escalation, workflow routing, stakeholder visualisation, exception handling,
Part-11 audit trail, P1--P9 emission.\\
\textit{Records: P1 (circuit fingerprint), P2 (compilation lineage), P6 (error-mitigation lineage), P7 (estimator/CI), P8 (cryptographic seal over P1--P7/P9).}
}}

\smallskip
\noindent\hspace{2em}$\Downarrow$ \textit{governed dispatch and lineage capture}
\smallskip

\noindent\colorbox{gray!15}{\parbox{\dimexpr\textwidth-2\fboxsep}{%
\textbf{Layer~2: hybrid-substrate layer}\\
Classical preprocessing/postprocessing, HPC baselines, AI rescoring, GPU services,
CUDA-Q/NVQLink, Ising-class or deterministic decoder lineage.\\
\textit{Records: P9 (ML-decoder lineage, sealed at this point).}
}}

\smallskip
\noindent\hspace{2em}$\Downarrow$ \textit{escalated kernel call when active-space and evidentiary criteria justify it}
\smallskip

\noindent\colorbox{gray!15}{\parbox{\dimexpr\textwidth-2\fboxsep}{%
\textbf{Layer~1: quantum-kernel layer}\\
Fermionic Hamiltonian, mapping and compilation, SqDRIFT sampling, backend execution
(IBM Heron~r3), shot distribution, estimator, confidence interval, mitigation metadata.\\
\textit{Records: P3 (device ID + calibration epoch), P4 (per-gate error budget snapshot), P5 (shot count + raw outcome distribution).}
}}

\medskip
\noindent\textbf{Governed output} $\leftarrow$
Observable estimate, uncertainty, escalation rationale, backend lineage,
P8 cryptographic seal over P1--P7/P9, reviewer-facing provenance package
\end{minipage}%
}
\caption{Three-layer QC--AI--HPC orchestration architecture. Each layer is annotated with
the P-records emitted at that execution stage. The quantum kernel is invoked selectively;
the hybrid substrate supplies surrounding AI/HPC services; the orchestration layer preserves
provenance and auditability across backend-specific execution paths.}
\label{fig:arch}
\end{figure}

\subsection{Layer 1: quantum-kernel layer}
\label{sec:layer1}

% CORRECTION #2 / #3: Animate which P-records are stamped at Layer 1 and their regulatory purpose.
Layer~1 houses the quantum kernel: QPU plus quantum algorithm.  In the dissertation proposal,
this layer is anchored by SqDRIFT on IBM Heron~r3.  SqDRIFT matters because the target
problem class is fermionic-Hamiltonian simulation with stochastic drift sampling.  IBM's
public materials describe a March 2026 quantum-computing contribution to a molecular discovery
published in \textit{Science}; secondary reporting describes the use of SqDRIFT on 100~qubits
to explore a 32-electron active space beyond an approximately 18-electron exact-classical
ceiling~\cite{ibm2026molecule,qcreport2026}.

At the moment of backend dispatch, Layer~1 stamps three provenance records.  The system emits
\textbf{P3}---device identifier and calibration epoch timestamp---to satisfy the FDA Part~11
requirement that the specific computing resource and its current characterisation be traceable
in the audit log; without P3, a regulator cannot verify that the hardware was within a valid
calibration window when the run was authorised.  Simultaneously, the system captures
\textbf{P4}---the per-gate error budget snapshot (one- and two-qubit fidelities, readout-error
matrices, T1/T2 values)---to provide the quantitative hardware state against which the
output's uncertainty bound can later be re-derived.  After shot execution, the system stores
\textbf{P5}---the full bitstring-level shot histogram---enabling independent re-derivation
of any aggregated estimator from the raw data, which is the probabilistic analogue of raw-data
preservation required under 21~CFR~Part~11~\S{}11.10(e).

The kernel layer is deliberately narrow.  It executes the escalated electronic-structure
subproblem under a recorded circuit, backend, calibration epoch, shot count, mitigation
policy, estimator, and decoder lineage.  That narrowness protects the paper from overstating
quantum results and clarifies the systems contribution.

\subsection{Layer 2: hybrid-substrate layer}
\label{sec:layer2}

% CORRECTION #2 / #3: Animate P9 at Layer 2 and its regulatory purpose.
Layer~2 contains the classical and AI substrate around the kernel.  It includes HPC
preprocessing, DFT/DMRG-style baselines where appropriate, conformational filtering, GPU
acceleration, AI-assisted rescoring, calibration support, and QEC decoding.  NVIDIA's April
2026 Ising announcement describes open quantum AI models for calibration and QEC with claimed
decoding speed and accuracy improvements over traditional approaches, and states that Ising
complements CUDA-Q and integrates with NVQLink for real-time control and
QEC~\cite{ising2026nvidia}.  NVIDIA's CUDA-Q QEC 0.6 documentation further describes
\textit{cudaq-realtime}, NVQLink-connected callbacks between GPUs and quantum controllers,
and real-time decoder pipelines including Ising CNN pre-decoders paired with
PyMatching~\cite{cudaqec2026}.

When an AI-based QEC decoder is active in Layer~2, the system emits \textbf{P9}---ML-decoder
lineage: decoder identifier and version, weights checksum, training-corpus manifest hash,
inference-time configuration, and execution environment.  This record exists to satisfy the
FDA Good Machine Learning Practice principle of model lifecycle control and
transparency~\cite{fda2021gmlp}: if a learned decoder shapes the inferred correction path, its
identity and configuration become part of the regulated computational record, not a hidden
implementation detail.  The cryptographic seal in P8 anchors P9 together with P1--P7, making
any post-hoc modification of decoder lineage detectable.  For deterministic decoders, P9
records the algorithm name and software version, preserving the same auditability at lower
lineage complexity.

\subsection{Layer 3: governed orchestration and provenance layer}
\label{sec:layer3}

% CORRECTION #2 / #3: Animate P1, P2, P6, P7, P8 at Layer 3 with regulatory rationale.
Layer~3 is the information-systems contribution.  It coordinates the first two layers and
emits the governance artifacts.  It implements selective escalation, records backend and
compilation choices, manages exceptions, presents stakeholder views, and produces the P1--P9
provenance record.

Specifically, Layer~3 stamps the following records in sequence.  At circuit compilation,
\textbf{P1}---the cryptographic hash of the post-transpilation compiled circuit---is
recorded; this is the regulatory anchor that makes the executed computation uniquely
identifiable and tamper-evident, fulfilling the audit-trail integrity requirement of
21~CFR~Part~11~\S{}11.10(a).  At fermion-to-qubit mapping, \textbf{P2}---compilation
lineage including source Hamiltonian, mapping choice, Trotter order/step, and SqDRIFT
sampling parameters---is populated; this record satisfies the traceability requirement
that every data transformation be documented so the output can be reconstructed from
inputs~\cite{fda21cfr11,fda2003guidance}.  After quantum execution and mitigation,
\textbf{P6}---the complete error-mitigation lineage, logging every transformation applied
to raw outcomes, method versions, and parameters---is stored.  Layer~3 then computes
\textbf{P7}---the final statistical estimator and 95\% confidence interval bounded by
shot count and noise-model variance---and finally applies \textbf{P8}, the time-stamped
cryptographic seal over P1--P7 and P9, anchoring the entire record against subsequent
modification in compliance with the electronic-signature and record-locking provisions
of 21~CFR~Part~11.

Layer~3 also makes the vendor split manageable.  If the workflow runs on IBM Heron through
Qiskit Runtime, compiler and calibration fields are populated for that path.  If a later
workflow uses a CUDA-Q/NVQLink-enabled QPU with an Ising-class decoder, the substrate and
P9 fields change while the workflow contract remains intact.  The architecture is therefore
portable without pretending that all backend behaviour is interchangeable.

%% ── §5 Selective Escalation ──────────────────────────────────────────────────
\section{Selective Computational Escalation}

The system does not ``use quantum computing everywhere.''  AI handles pattern detection,
rescoring, and interpretive support where such models are appropriate.  HPC handles broad
preprocessing, classical simulation, and postprocessing where classical methods remain
tractable and reproducible.  Quantum simulation is invoked only where active-space size,
correlation structure, and accuracy requirements justify an escalated kernel.

The BQP-relevant classification is therefore operational, not absolutist.  It relies on
the established quantum-computing framing that selected local-Hamiltonian and
molecular-energy estimation problems admit quantum
formulations~\cite{abrams1999,aspuru2005,bennett1997}.  If a classical method later demonstrates
chemical accuracy for the specified active-space boundary within tractable resources, the
orchestration layer routes the workload to that classical backend.  This revisability is
a feature: the architecture governs escalation rather than advocating one substrate
regardless of evidence.

The BQP-relevant classification operationalises as a three-condition escalation rule.
A workload escalates to Layer~1 if and only if all three conditions in Table~\ref{tab:escalation}
are satisfied simultaneously; otherwise it routes to Layer~2 or remains classical, with the
routing rationale recorded in the P2 provenance field.

\begin{table}[H]
\centering
\caption{Three-condition escalation decision rule with C275F local-site instantiation.}
\label{tab:escalation}
\small
\setlength{\tabcolsep}{6pt}
\begin{tabular}{@{}
  p{2.6cm}
  p{5.4cm}
  p{4.8cm}
  >{\centering\arraybackslash}p{0.8cm}
  @{}}
\toprule
\textbf{Condition} & \textbf{Criterion} & \textbf{C275F local-site value} & \textbf{Met?} \\
\midrule
\textbf{C1: Active-space size} &
  Active electrons exceed classical FCI ceiling (${\sim}$18e$^{-}$ under tractable resources) &
  ${\sim}$24 active electrons at local binding site &
  $\checkmark$ \\[4pt]
\textbf{C2: Algorithmic formulation} &
  Fermionic Hamiltonian admits a known quantum algorithmic formulation &
  Jordan--Wigner mapping available; SqDRIFT applicable &
  $\checkmark$ \\[4pt]
\textbf{C3: Hardware readiness} &
  Calibration epoch meets P3 timestamp policy for target backend &
  IBM Heron~r3 calibration epoch current at dispatch &
  $\checkmark$ \\
\midrule
\multicolumn{4}{@{}p{13.8cm}@{}}{%
  \textbf{Decision:} All three conditions satisfied $\Rightarrow$ Route to
  Layer~1 (quantum kernel). Any condition fails $\Rightarrow$ Route to
  Layer~2 (hybrid substrate); record routing rationale in P2.} \\
\bottomrule
\end{tabular}
\end{table}

\textbf{Distinction from quantum-enhanced classical workflows.}  The BQP-relevant
classification operationalises a boundary that existing quantum drug-discovery literature
does not always make explicit.  Vakili et al.\ applied a hybrid quantum--classical
generative model---a quantum circuit Born machine (QCBM) on a 16-qubit IBM processor
paired with an LSTM network---to KRAS inhibitor design, yielding experimentally validated
hits and a 21.5\% improvement in candidate quality over the classical
baseline~\cite{vakili2025kras}.  This is quantum-\emph{enhanced} drug discovery: KRAS
has resolved crystal structures, a substantial library of known inhibitors, and classical
algorithms already generate viable candidates.  Quantum computing refines the search;
it does not replace the only available computational path.

The mutations targeted by this dissertation occupy a categorically different regime.
TP53~C275F, STK11~F354L, STK11~D194N, KEAP1~G333C, and KEAP1~R320Q have no
binding pocket accessible by classical structure-based methods, no prior computational
characterisation at chemical accuracy, and active-space sizes at which DFT introduces
systematic errors of 2--4~kcal/mol---sufficient to invert the sign of a binding
prediction.  Quantum simulation is therefore not an enhancement but a
\emph{necessity}: it is the only computational path to a physically grounded
$\Delta E$ estimate.  The three-condition escalation rule in Table~\ref{tab:escalation}
encodes this distinction as a routing decision; the Vakili et al.\ KRAS workflow would
not satisfy C1 (active-space ceiling not exceeded by classical methods) and would
route to Layer~2 rather than the quantum kernel.

%% ── §6 Provenance Schema and Lifecycle ────────────────────────────────────────
\section{Part-11-Governed Provenance for Probabilistic Quantum Simulation}

FDA~21~CFR~Part~11 governs electronic records and electronic signatures in FDA-regulated
settings~\cite{fda21cfr11,fda2003guidance}.  Quantum simulation complicates record governance
because a valid output is normally a distribution of measurement outcomes, not a single
deterministic value.  A rerun may be statistically equivalent without being bitwise identical.
The architecture therefore operationalises reproducibility as statistical equivalence under
recorded circuit, backend, calibration, error-mitigation, estimator, and decoder conditions.

Table~\ref{tab:provenance} defines the nine-element provenance record.  The wording follows
the dissertation proposal because this schema is the paper's core governance contribution.

\begin{table}[H]
\centering
\caption{P1--P9 provenance schema for governed probabilistic quantum simulation.}
\label{tab:provenance}
\small
\begin{tabularx}{\textwidth}{@{}lX@{}}
\toprule
\textbf{Field} & \textbf{Record content} \\
\midrule
\textbf{P1} & Circuit fingerprint: cryptographic hash of the post-transpilation compiled circuit. \\
\textbf{P2} & Compilation lineage: source Hamiltonian, fermion-to-qubit mapping, Trotter order/step, SqDRIFT sampling parameters. \\
\textbf{P3} & Device identifier and calibration epoch: backend, processor revision, calibration timestamp. \\
\textbf{P4} & Per-gate error budget snapshot: one- and two-qubit fidelities, readout-error matrices, T1/T2 at the calibration epoch. \\
\textbf{P5} & Shot count and raw outcome distribution: full bitstring-level histogram permitting independent re-derivation of any aggregated estimator. \\
\textbf{P6} & Error-mitigation lineage: every transformation applied to raw outcomes, including zero-noise extrapolation, probabilistic error cancellation, twirling, parameters, and method versions. \\
\textbf{P7} & Statistical estimator and confidence interval: final reported observable, estimator function, 95\% CI bounded by shot-count and noise-model variance. \\
\textbf{P8} & Cryptographic seal: time-stamped digital signature over P1--P7 and P9, anchoring the record against subsequent modification. \\
\textbf{P9} & ML-decoder lineage: when QEC inference uses a learned model---decoder identifier and version, weights checksum, training-corpus manifest hash, inference-time configuration (precision, batch size, decoder threshold), and execution environment (GPU class, CUDA/runtime version).  For deterministic decoders: algorithm name and software version. \\
\bottomrule
\end{tabularx}
\end{table}

P9 is not an optional flourish.  It is what makes the architecture forward-compatible with
learned QEC and error-mitigation workflows.  FDA's good machine learning practice principles
emphasise model lifecycle control and transparency for learned
systems~\cite{fda2021gmlp}.  NVIDIA's Ising/CUDA-Q materials show why the same logic becomes
relevant inside a quantum pipeline when learned decoders become part of
execution~\cite{ising2026nvidia,cudaqec2026}.

\subsection*{Provenance in motion: P1--P9 lifecycle for a C275F local-site run}
\label{sec:provenance-lifecycle}

% CORRECTION #6: Step-by-step numbered walkthrough annotating each P-record
% with the layer that emits it and the regulatory purpose it serves.
The following walk-through traces a C275F local binding-site run through the nine-step
provenance lifecycle.  Each step is annotated with the P-record emitted and the
Part-11/GMLP regulatory need it satisfies.

\begin{enumerate}[leftmargin=2em,itemsep=4pt]

\item \textbf{Circuit compilation (Layer~3).}\;
After the Jordan--Wigner-mapped fermionic Hamiltonian is transpiled to a Heron-native
gate set, Layer~3 computes a SHA-256 hash of the compiled circuit and writes it as \textbf{P1}.
\textit{Regulatory purpose:} satisfies 21~CFR~Part~11~\S{}11.10(a) audit-trail integrity---any
subsequent modification of the compiled circuit would invalidate P1 and be detected on audit.

\item \textbf{Compilation lineage capture (Layer~3).}\;
The system writes \textbf{P2}: source Hamiltonian identity, choice of Jordan--Wigner
mapping, Trotter order and step size, and SqDRIFT sampling configuration.
\textit{Regulatory purpose:} satisfies the traceability requirement that every
data-transformation parameter be documented so that the output can be reconstructed from
inputs, as required under 21~CFR~Part~11~\S{}11.10(b).

\item \textbf{Backend dispatch (Layer~1).}\;
At the moment of IBM Heron~r3 job submission, Layer~1 stamps \textbf{P3}: device
identifier, processor revision, and calibration epoch timestamp.
\textit{Regulatory purpose:} provides the regulatory anchor that the specific computing
resource was characterised and within a valid calibration window when the run was
authorised; essential for audit reconstruction.

\item \textbf{Hardware state snapshot (Layer~1).}\;
Concurrently with dispatch, Layer~1 stores \textbf{P4}: per-gate one- and two-qubit
fidelities, readout-error matrices, and T1/T2 decoherence times at the calibration epoch.
\textit{Regulatory purpose:} the quantitative hardware state against which the output
uncertainty bound can be independently re-derived; supports the reproducibility standard
for probabilistic outputs.

\item \textbf{Shot execution and histogram storage (Layer~1).}\;
After quantum execution, Layer~1 stores \textbf{P5}: the full bitstring-level shot
histogram.
\textit{Regulatory purpose:} equivalent to raw-data preservation under 21~CFR~Part~11
\S{}11.10(e); enables independent re-derivation of any aggregated estimator from raw data.

\item \textbf{Decoder lineage (Layer~2---conditional).}\;
If an Ising-class or other learned QEC decoder is active, Layer~2 emits \textbf{P9}: decoder
identifier, weights checksum, training-corpus manifest hash, inference-time configuration,
and execution environment.  For deterministic decoders, P9 records algorithm name and
software version.
\textit{Regulatory purpose:} satisfies FDA GMLP model lifecycle control requirements;
makes decoder identity part of the regulated computation record, not a hidden dependency.

\item \textbf{Error-mitigation lineage (Layer~3).}\;
Layer~3 writes \textbf{P6}: a complete ordered log of every mitigation transformation
applied to raw outcomes---zero-noise extrapolation factors, probabilistic error-cancellation
quasi-probabilities, twirling parameters---along with method identifiers and software
versions.
\textit{Regulatory purpose:} enables independent auditors to reproduce the post-processed
estimate from the raw P5 histogram.

\item \textbf{Statistical estimator and confidence interval (Layer~3).}\;
Layer~3 computes and stores \textbf{P7}: the final reported observable, the estimator
function applied, and the 95\% confidence interval bounded by shot count and noise-model
variance derived from P4.
\textit{Regulatory purpose:} provides the primary reported result in a form that quantifies
and bounds uncertainty, replacing the na\"{\i}ve point-value convention inappropriate for
probabilistic quantum outputs.

\item \textbf{Cryptographic seal (Layer~3).}\;
Layer~3 applies \textbf{P8}: a time-stamped digital signature over the concatenated
P1--P7 and P9 record.
\textit{Regulatory purpose:} satisfies 21~CFR~Part~11 electronic-signature and
record-locking provisions; any post-hoc modification of any P-field is detectable by
re-verifying P8.

\end{enumerate}

The result is a self-contained, tamper-evident provenance package that travels with
the output and can be presented to a regulatory reviewer, an independent replicator, or a
stakeholder audit committee without requiring reconstruction from logs.

\subsection{Access Control and User Attribution as Part-11 Infrastructure}
\label{sec:rbac-jwt}

FDA~21~CFR~Part~11 \S11.10(d) requires that system access be limited to authorised
individuals.  \S11.10(i) requires that audit trails capture the identity of the person
performing each recorded action.  The current implementation satisfies both requirements
through two IS-level mechanisms that operate within the governed orchestration layer
(Layer~3) rather than the quantum or hybrid substrate layers.

\textbf{Role-based access control (RBAC).}  On authentication, the platform queries a
persistent user-profile store for the authenticated user's assigned role.  Administrative
functions---including access to all simulation runs across users, the provenance audit log,
and user management---are gated behind a role check and invisible to non-administrative
users.  The access boundary is enforced by a database-resident role assignment, not by
convention or self-restraint, mapping directly to the \S11.10(d) requirement that the
system boundary be actively enforced.

\textbf{JWT-based user attribution.}  Every simulation run carries a JSON Web Token (JWT)
credential.  The platform fetches the authenticated session JWT and transmits it as a
Bearer credential with each simulation request.  The backend decodes the JWT to extract
the user identifier and records it against the provenance row in the database, creating an
unbroken chain from authenticated login to sealed P1--P9 record.  This satisfies
\S11.10(i) without requiring the backend to maintain independent credential state---the
identity provider handles verification separately.

Together, RBAC and JWT attribution are IS design decisions with direct regulatory
consequences.  A quantum simulation pipeline that emits complete, SHA-256-sealed P1--P9
records but cannot identify which authorised user initiated each run is not
Part-11-compatible.  These mechanisms extend the governance contract of Layer~3 into the
access and identity layer and are accordingly treated as part of the IS contribution of
this paper.

%% ── §7 Research Artifact and Prototype ───────────────────────────────────────
\section{Research Artifact and Prototype}

The research artifact is a QC--AI--HPC orchestration platform and associated governance
schema.  The public prototype associated with the dissertation proposal demonstrates the
intended stakeholder workflow, including a multi-screen platform, architecture visualisation,
and molecular-viewer interface~\cite{cohen2026prototype}.  In this article, the prototype is
treated as an executable design specification rather than as evidence of quantum-computed
binding energies.

Phase~3A of the prototype executes on a PennyLane noiseless simulator backend using a minimal
two-electron active space with physically motivated proxy Hamiltonian coefficients---manually
varied Jordan--Wigner terms that serve as stand-ins for the full PySCF-derived molecular
Hamiltonians targeted in Phase~3B\@.  This is deliberate architectural scaffolding, not a
scientific claim: the proxy Hamiltonians are sufficient to exercise every pipeline stage---NGS
ingestion, Hamiltonian construction, VQE execution, provenance assembly, cryptographic sealing,
and database storage---while Phase~3B replaces them with quantum chemistry computations on the
RCSB crystal structures already resolved by the platform's structure-source layer.  The five
NSCLC mutation targets produce distinct proxy coefficients, distinct ground-state energies, and
distinct P8 SHA-256 seals, demonstrating that the provenance schema correctly distinguishes
between simulation runs before molecular-accuracy Hamiltonians are introduced in Phase~3B\@.

This distinction is necessary for a defensible arXiv submission.  The article makes workflow,
record, and evaluation criteria inspectable before full physical-hardware execution exists.
The architecture can therefore be reviewed by information-systems scholars for artifact design,
by quantum specialists for backend and provenance assumptions, by biomedical readers for
translational boundaries, and by regulatory readers for record-governance logic.

%% ── §8 Evaluation Protocol ────────────────────────────────────────────────────
\section{Evaluation Protocol}

The evaluation is triangulated rather than single-metric.  The dissertation proposal specifies
four layers, which are adapted here into an arXiv-ready validation plan following the FEDS
framework~\cite{venable2016}.

\subsection{Layer 1: technical evaluation}

% CORRECTIONS #4 / #7: Each metric presented as a stress test derived from a real constraint;
% failure-mode sentence added for each threshold.
Layer~1 measures the artifact itself through automated harnesses.  Each acceptance criterion
below is derived from a hardware, workflow, or regulatory constraint; the failure mode if
the threshold is not met is stated explicitly.

\begin{itemize}[leftmargin=1.5em,itemsep=4pt]

\item \textbf{End-to-end completion $\geq 95\%$ over simulated runs.}\;
The 95\% threshold is set at the point where the remaining 5\% failure rate matches the
expected noise floor of simulator-level qubit drop-out and decoherence under current
IBM Heron~r3 error rates; requiring 100\% would be unrealistic for stochastic quantum
workflows, while tolerating more than 5\% failures would indicate a systematic fault in
orchestration logic rather than hardware noise.
\textit{Failure mode:} completion below 95\% signals a broken escalation rule or an
unhandled exception path, rendering the architecture non-deployable in a regulated
oncology pipeline.

\item \textbf{100\% P1--P9 provenance emission for all recorded outputs.}\;
Derived from 21~CFR~Part~11's requirement that every regulated computational output carry
a complete, unambiguous audit record; partial records cannot satisfy regulatory review.
\textit{Failure mode:} any missing P-field in a regulated run constitutes an audit gap
that would disqualify the output from a submission package under Part~11 conditions.

\item \textbf{Cryptographic-seal (P8) verification latency $\leq 5\%$ of total workflow
wall-clock time.}\;
This threshold is grounded in the real-time feedback loop between GPU (Layer~2) and QPU
(Layer~1): NVIDIA's NVQLink QEC pipeline operates within qubit coherence windows on the
order of microseconds to milliseconds~\cite{cudaqec2026}.  If cryptographic overhead
exceeds 5\% of wall-clock time, the seal operation competes with the real-time
error-correction callback window, breaking the in-loop correction guarantee before qubits
decohere.
\textit{Failure mode:} P8 latency above 5\% causes the architecture to lose regulatory
viability in time-sensitive oncology pipelines where real-time QEC is required.

\item \textbf{100\% audit-trail completeness across the local-to-full active-space
boundary.}\;
Derived from the regulatory requirement that the same record structure applies regardless
of circuit depth or qubit count; an architecture that loses provenance at scale is not
production-ready.
\textit{Failure mode:} provenance gaps at the full ($\sim$88-qubit) tier reveal a scaling
fault in the orchestration layer that would prevent frontier-scale regulatory use.

\item \textbf{$\geq 90\%$ detection of injected fault scenarios.}\;
The 90\% threshold reflects the realistic bound on automated fault detection given the
stochastic nature of quantum outputs; 100\% is unachievable without access to a ground
truth that does not exist for probabilistic hardware.  Fault injection follows the taxonomy
used in the case-based walkthrough (Section~\ref{sec:walkthrough}).
\textit{Failure mode:} detection below 90\% implies systematic blind spots in the
exception-handling layer, undermining the governance contract.

\end{itemize}

\subsection{Layer 2: comparative evaluation}

Layer~2 compares the architecture against four narrower alternatives: direct Qiskit Runtime
execution, Orquestra-style orchestration, CUDA-Q/NVQLink substrate pathways, and a manual
notebook workflow.  The comparison is not meant to prove that the proposed system is faster
at every computational step.  It tests whether the alternatives provide the same combined
capabilities: oncology-mutation-specific construction, selective escalation, AI/HPC
integration, stakeholder visualisation, statistical-equivalence validation, and
Part-11-grade provenance.

% CORRECTION #4 / #7: Failure modes and physics/social rationale added.
The proposed architecture passes the comparative evaluation if and only if: (a)~0\% of
simulated runs produce incomplete P1--P9 records, where all four alternatives are expected
to produce partial or absent audit trails---\textit{failure mode: any incomplete record
disqualifies the architecture for FDA-regulated submissions}; (b)~Part-11-grade
cryptographic seal (P8) verification succeeds with latency overhead $\leq 5\%$ of total
workflow wall-clock time, justified by the NVQLink real-time coherence constraint described
above---\textit{failure mode: P8 latency above 5\% breaks the real-time QEC callback loop};
and (c)~stakeholder visualisation is available across all three active-space tiers defined
in Table~\ref{tab:tiers}, a capability absent from all four alternatives---\textit{failure
mode: missing tier coverage implies the governance contract does not extend to
frontier-scale runs, limiting translational utility}.

\subsection{Layer 3: expert-panel evaluation}

% CORRECTION #4 / #7: Social/epistemic rationale for scoring threshold added.
Layer~3 uses expert review to test cross-disciplinary validity.  The dissertation proposal
specifies a panel of four to six reviewers spanning quantum chemistry/quantum computing,
oncology/molecular pathology, pharma R\&D/computational drug discovery,
information-systems/DSR methodology, and regulatory/clinical-informatics expertise.  Review
criteria map to the Hevner DSR guidelines: artifact existence, problem relevance, design
evaluation, contribution, rigor, design-search process, and communication~\cite{hevner2004}.

Reviewers score the artifact against each of Hevner's seven DSR criteria on a 1--5 integer
scale.  The expert-panel evaluation passes if the mean score across all reviewers reaches at
least \textbf{3.5} on every dimension and no criterion receives a majority of scores of 2 or
below.  The 3.5 threshold is meaningfully demanding in this setting because it represents
cross-domain agreement among oncologists, quantum chemists, information-systems scholars, and
regulatory specialists who operate under fundamentally different evidentiary standards and
whose consensus is difficult to achieve by coincidence; a mean below 3.5 would indicate that
the artifact fails to bridge at least one disciplinary boundary, reducing its translational
utility.  Where scores diverge by more than 2 points on any criterion, the panel convenes a
structured reconciliation discussion before a final score is recorded.
\textit{Failure mode:} a mean below 3.5 on any criterion signals that the artifact is
insufficiently legible or credible to one expert community, which would undermine adoption
in the multi-stakeholder oncology pathway.

\subsection{Layer 4: case-based walkthrough}
\label{sec:walkthrough}

Layer~4 traces a de-identified C275F scenario and a second mutation scenario through the
workflow.  The purpose is to show how the system works in context: where data enter, where
routing decisions occur, where the audit layer records or rejects events, which stakeholder
view is produced, and how anomalies are handled.  Fault scenarios injected include:
calibration-epoch staleness triggering a P3 policy violation; active-space boundary
re-evaluation demoting a run from Layer~1 to Layer~2; and a learned-decoder weight
mismatch producing a P9 integrity failure.  This layer complements the technical benchmarks
by testing whether the artifact is understandable and reviewable by its intended users.

%% ── §9 Discussion ────────────────────────────────────────────────────────────
\section{Discussion}

\subsection{Information-systems contribution}

The primary contribution is an architecture for governed computational escalation that
turns a collection of specialised tools into a reviewable information system.  Quantum
simulation, AI rescoring, GPU-accelerated decoding, and HPC preprocessing each produce
partial evidence under different assumptions; the orchestration layer makes those assumptions
explicit enough for audit, comparison, and stakeholder review.

The architecture's governance contract is defined by three domain-independent parameters:
active-space size relative to the classical tractability ceiling, availability of a known
quantum algorithmic formulation, and the compliance tier of the regulated workflow.  None
of these parameters are specific to TP53 or NSCLC.  Any fermionic simulation workload
whose active-space boundary falls at the classical--quantum frontier and whose regulatory
context requires an auditable escalation record is a candidate for the same three-layer
orchestration pattern.  Extension to other mutation classes, cancer types, or non-oncology
molecular simulation domains requires only reconfiguration of the active-space routing
policy and the stakeholder-facing visualisation layer.

\subsection{Emerging-computing contribution}

The emerging-computing contribution is bifurcation awareness.  The architecture recognises
that the field is splitting between quantum-kernel platforms and hybrid AI/HPC substrates.
IBM's Heron/Qiskit path and NVIDIA's CUDA-Q/NVQLink/Ising path are not forced into one
universal abstraction.  Instead, the workflow contract records which path is used and keeps
substrate-specific evidence available for review.

\subsection{Regulatory contribution}

The regulatory contribution is the reinterpretation of Part-11-style records for
probabilistic quantum outputs.  Bitwise replay is not the right standard for a shot-based
quantum simulation whose legitimate result is a distribution.  The proposed record
preserves the conditions under which statistical equivalence can be assessed.  P9 extends
this logic to learned decoders, where model lineage becomes part of the evidence chain.

%% ── §10 Limitations and Future Work ──────────────────────────────────────────
\section{Limitations and Future Work}

% CORRECTION #1: All clinical/chemistry disclaimers consolidated here only.
This paper does not report a completed full active-site C275F quantum simulation on physical
IBM Heron~r3 hardware and does not present binding-energy values as executed quantum
results.  It does not claim an FDA submission, a clinical-trial result, or a therapeutic
recommendation, and does not claim a formal classical-versus-quantum separation for C275F.
It does not empirically validate NVIDIA Ising-class decoder provenance in a live regulated
workflow; that is a future phase.

A May 2026 milestone by Merz et al.\ demonstrates that the hardware scale required for
Phase~3B is no longer speculative: 94 qubits were used on IBM Heron~r2 processors to
compute the electronic structure of a 12,635-atom trypsin–benzamidine complex via
EWF-TrimSQD at chemical accuracy~\cite{merz2026protein}.  C275F's full binding-site
target ($\sim$44 electrons, $\sim$88 qubits) falls comfortably within this demonstrated
range.  The architectural contribution of this dissertation is therefore validated against
hardware that can execute Phase~3B; the remaining step is institutional IBM Quantum Network
access, not algorithmic or hardware feasibility.

The architecture is evaluated initially on one mutation instance and one quantum-modality
anchor.  Generalisation is therefore analytical rather than statistical.  Multi-patient
analysis, multi-modality empirical validation across trapped-ion or neutral-atom systems,
and live cross-stack validation across IBM-native and NVIDIA-substrate pathways remain
future work.  These boundaries define the difference between a publishable architecture
paper and an overextended chemistry or clinical paper; they do not undermine the
information-systems claim.

%% ── §11 Conclusion ────────────────────────────────────────────────────────────
\section{Conclusion}

This work is strongest when read as an information-systems contribution: the problem is not
merely to compute a molecular number, but to design the system that makes such computation
governable.  The proposed QC--AI--HPC architecture organises
the work into a quantum-kernel layer, a hybrid-substrate layer, and a governed
orchestration/provenance layer.  The P1--P9 provenance schema animates that architecture,
stamping records at each execution stage with a clear regulatory purpose: P3, P4, and P5
anchor hardware-state traceability at Layer~1; P9 captures decoder lineage at Layer~2;
and P1, P2, P6, P7, and P8 seal the complete computational chain at Layer~3.  Together,
these records turn a probabilistic quantum pipeline into a governable, auditable
information system fit for regulated oncology settings.

%% ── License ──────────────────────────────────────────────────────────────────
\section*{License and arXiv Submission Note}

This manuscript is prepared for arXiv submission under the author-selected CC BY-NC-ND~4.0
license.  arXiv lists CC BY-NC-ND~4.0 as allowing copying and distribution of the material
in unadapted form for noncommercial purposes with attribution, and notes that the license
chosen for each version is irrevocable~\cite{arxiv2026license}.  The submitter should verify
compatibility with university, funder, and later journal policies before upload.

%% ── Acknowledgments ──────────────────────────────────────────────────────────
\section*{Acknowledgments}

This manuscript was prepared from the author's dissertation proposal and associated prototype
materials under the supervision of Prof.\ Itamar Shabtai.  The author remains responsible
for verifying final authorship metadata, institutional approvals, patient-data handling,
and arXiv submission settings before upload.

%% ── References ───────────────────────────────────────────────────────────────
\begin{thebibliography}{19}

\bibitem{cohen2026proposal}
Doron Cohen.
A Translational QC--AI--HPC Orchestration Platform for Non-druggable NSCLC Tumor-Suppressor Mutations.
Dissertation proposal website, 2026.
\url{https://nsclc-quantum-viewer.netlify.app/dissertation_revised}.
Accessed May~11, 2026.

\bibitem{abrams1999}
Daniel S. Abrams and Seth Lloyd.
Quantum algorithm providing exponential speed increase for finding eigenvalues and eigenvectors.
\textit{Physical Review Letters}, 83(24):5162--5165, 1999.
\doi{10.1103/PhysRevLett.83.5162}.

\bibitem{aspuru2005}
Al\'{a}n Aspuru-Guzik, Anthony~D. Dutoi, Peter~J. Love, and Martin Head-Gordon.
Simulated quantum computation of molecular energies.
\textit{Science}, 309(5741):1704--1707, 2005.
\doi{10.1126/science.1113479}.

\bibitem{heron2026ibm}
IBM Quantum.
Processor Types.
IBM Quantum documentation, 2026.
\url{https://quantum.cloud.ibm.com/docs/en/guides/processor-types}.
Accessed May~11, 2026.

\bibitem{ising2026nvidia}
NVIDIA.
NVIDIA Launches Ising, the World's First Open AI Models to Accelerate the Path to Useful Quantum Computers.
NVIDIA investor press release, 2026.

\bibitem{cudaqec2026}
NVIDIA CUDA-Q.
CUDA-Q QEC 0.6 Enables Real-Time QEC with NVQLink.
NVIDIA CUDA-Q blog, 2026.
\url{https://nvidia.github.io/cuda-quantum/blogs/blog/2026/04/14/cudaq-qec-0.6/}.
Accessed May~11, 2026.

\bibitem{hevner2004}
Alan~R. Hevner, Salvatore~T. March, Jinsoo Park, and Sudha Ram.
Design science in information systems research.
\textit{MIS Quarterly}, 28(1):75--105, 2004.

\bibitem{peffers2007}
Ken Peffers, Tuure Tuunanen, Marcus~A. Rothenberger, and Samir Chatterjee.
A design science research methodology for information systems research.
\textit{Journal of Management Information Systems}, 24(3):45--77, 2007.
\doi{10.2753/MIS0742-1222240302}.

\bibitem{venable2016}
John Venable, Jan Pries-Heje, and Richard Baskerville.
FEDS: A Framework for Evaluation in Design Science Research.
\textit{European Journal of Information Systems}, 25(1):77--89, 2016.
\doi{10.1057/ejis.2014.36}.

\bibitem{ibm2026molecule}
IBM.
IBM and University Researchers Create a Never-Before-Seen Molecule and Prove its Exotic Nature with Quantum Computing.
IBM Newsroom, 2026.

\bibitem{qcreport2026}
Quantum Computing Report.
IBM and Researchers Synthesize First Half-M\"{o}bius Molecule Validated by Quantum Computing.
Quantum Computing Report, 2026.

\bibitem{bennett1997}
Charles~H. Bennett, Ethan Bernstein, Gilles Brassard, and Umesh Vazirani.
Strengths and weaknesses of quantum computing.
\textit{SIAM Journal on Computing}, 26(5):1510--1523, 1997.
\doi{10.1137/S0097539796300933}.

\bibitem{fda21cfr11}
U.S. Food and Drug Administration.
21~CFR Part~11: Electronic Records; Electronic Signatures.
Electronic Code of Federal Regulations, 2026.

\bibitem{fda2003guidance}
U.S. Food and Drug Administration.
Guidance for industry: Part~11, electronic records; electronic signatures---scope and application.
FDA guidance document, 2003.

\bibitem{fda2021gmlp}
U.S. Food and Drug Administration.
Good Machine Learning Practice for Medical Device Development: Guiding Principles.
FDA web resource, 2021.

\bibitem{cohen2026prototype}
Doron Cohen.
NSCLC Quantum Viewer and QC--AI--HPC Prototype.
Prototype website, 2026.
\url{https://nsclc-quantum-viewer.netlify.app/}.
Accessed May~11, 2026.

\bibitem{vakili2025kras}
Mohammad~Ghazi Vakili, Christoph Gorgulla, Jamie Snider, et~al.
Quantum-computing-enhanced algorithm unveils potential KRAS inhibitors.
\textit{Nature Biotechnology}, 43:1954--1959, 2025.
\doi{10.1038/s41587-024-02526-3}.

\bibitem{merz2026protein}
Kenneth~M. Merz et~al.
Crossing the 12,000-atom barrier with heterogeneous quantum-classical supercomputing:
quantum chemistry of protein-ligand complexes.
arXiv preprint arXiv:2605.01138, 2026.
Cleveland Clinic Foundation, RIKEN, and IBM Quantum.
\url{https://arxiv.org/abs/2605.01138}.

\bibitem{arxiv2026license}
Licenses.
arXiv help page, 2026.
\url{https://info.arxiv.org/help/license/index.html}.
Accessed May~11, 2026.

\end{thebibliography}

\end{document}
